\hypertarget{tcltk-and-gui-foundations-tkinter}{%
\chapter{\texorpdfstring{Tcl/Tk and GUI Foundations
(\texttt{tkinter})}{Tcl/Tk and GUI Foundations (tkinter)}}\label{tcltk-and-gui-foundations-tkinter}}

\passthrough{\lstinline!tkinter!} is the standard Python interface to
the Tk GUI toolkit.

\hypertarget{the-c-bridge-_tkinter}{%
\subsection{\texorpdfstring{59.1 The C-Bridge:
\texttt{\_tkinter}}{59.1 The C-Bridge: \_tkinter}}\label{the-c-bridge-_tkinter}}

\passthrough{\lstinline!tkinter!} is not written in Python. It is a
wrapper around the \textbf{Tcl/Tk} C library. * \textbf{The Tcl
Interpreter}: When you instantiate \passthrough{\lstinline!Tk()!}, a
full Tcl interpreter is created inside your Python process. *
\textbf{Command Marshalling}: When you call
\passthrough{\lstinline!button.configure(text="Click")!}, Python
marshals the arguments into Tcl strings and executes them in the Tcl VM.

\hypertarget{the-main-loop-and-event-concurrency}{%
\subsection{59.2 The Main Loop and Event
Concurrency}\label{the-main-loop-and-event-concurrency}}

GUIs are event-driven. \passthrough{\lstinline!root.mainloop()!} enters
a blocking loop that waits for OS events (mouse clicks, key presses). *
\textbf{Thread Safety}: Tk is not thread-safe. All GUI updates must
happen on the main thread. * \textbf{\passthrough{\lstinline!after()!}}:
Use \passthrough{\lstinline!root.after(ms, callback)!} to schedule
Python functions without blocking the GUI event loop. This is
effectively a simple cooperative multitasking scheduler built on top of
the Tk event queue.

\begin{center}\rule{0.5\linewidth}{0.5pt}\end{center}
